\section{诊断与稳健性}
\label{sec:diagnostics}

界面中位数对稀疏受体几乎没有区分度。9000个原发灶配体受体--切片组合里，界面中位数大于0的占6.1\%。1000对里，9例中位数等于0的占95.3\%。C3在成纤维高分区的检出比例可以到0.50--0.96，C3AR1在界面上多为0.14--0.36，乘积仍有超过一半的界面点位为0，中位数因此落在0。GC4是例外：C3--C3AR1的界面中位数为0.912，空分布第95百分位为0.845，置换$p$为0.005。把9例再取中位数时，这一例被其余8个0盖住。

留一规则比“刚好5例”更紧。APP--CD74的5例过线分布在部分切片上，去掉GC1、GC4、GC6、GC8或GC9中的任意一张，过线数就少于5。CD99--CD99和SEMA4D--PLXNB2同样卡在5例。置换次数为200，第95百分位对应的切片内置换$p$约为0.05；中位数置换$p$为0.030，表示至少一半切片的置换$p$不高于这个量级，另一半可以接近1。因此5/9过线不等于9例一致。

接触比例的空分布集中在0.95--0.97，观察值集中在0.81--0.94。两组都高，差值的方向在10张切片上一致：观察值更低。这个方向说明当前半径和前25\%阈值下，真实标签的接触不多于随机标签。
